%%%%%%%%%%%%%%%%%%%%%%%%%%%%%%%%%%%%%%%%%%%%%%%%%%%%%%%%%%%%%%%%%%%%%%
% Lớp 10 | Đại số tổ hợp | Hoán vị. Chỉnh hợp. Tổ hợp
% Dạng 0D8?2-2: Bài toán có biểu thức 𝑃, 𝐶, 𝐴
% Số câu: 28  (sinh tự động bởi tach_cau.py; từ nay thầy sửa trực tiếp file này)
%%%%%%%%%%%%%%%%%%%%%%%%%%%%%%%%%%%%%%%%%%%%%%%%%%%%%%%%%%%%%%%%%%%%%%

%%%%%%%%%%%%%%%% TRẮC NGHIỆM 4 PHƯƠNG ÁN %%%%%%%%%%%%%%%%

% nguồn: MaHoaNH_TracNghiem_Lop10
\begin{ex}%[0D8N2-2]%[ma:186cc81d]
	Cho tập $M$ có $10$ phần tử. Số tập con gồm $2$ phần tử  của $M$ là
	 \choice
		 {$\mathrm{A}_{10}^8$}
		{$\mathrm{A}_{10}^2$} 
		{\True $\mathrm{C}_{10}^2$} 
		{$10^2$} 
	 \loigiai{
	Số tập con gồm $2$  phần tử của tập $M$ có $10$ phần tử là $\mathrm{C}_{10}^2$.}
\end{ex}

% nguồn: MaHoaNH_TracNghiem_Lop10
\begin{ex}%[0D8N2-2]%[ma:8f02be3c]
	 Trong mặt phẳng cho $6$ điểm phân biệt. Có bao nhiêu véc-tơ khác $\overrightarrow{0}$ có điểm đầu và điểm cuối là $2$ trong $6$ điểm đã cho?
	 \choice
		 {$1440$}
		 {\True $30$}
		{$12$} 
		 {$15$}
	\loigiai{Số véc-tơ khác  $\overrightarrow{0}$ có điểm đầu và điểm cuối là $2$ trong $6$ điểm đã cho là $\mathrm{A}_6^2=30$.
	 }
\end{ex}

% nguồn: MaHoaNH_TracNghiem_Lop10
\begin{ex}%[0D8N2-2]%[ma:a7f9fcca]
	Số cách sắp xếp $6$ nam sinh và $4$ nữ sinh vào một dãy ghế hàng ngang có  $10$ chỗ ngồi là
	\choice 
		{$6!\cdot 4!$} 
		 {\True $10!$}
		{$6!+4!$} 
		 {$6!-4!$}
	\loigiai{ 
		Có $10$ học sinh xếp vào một dãy ngế hàng ngang có $10$ chỗ ngồi nên có  $10!$ cách.
	} 
\end{ex}

% nguồn: MaHoaNH_TracNghiem_Lop10
\begin{ex}%[0D8N2-2]%[ma:1ef6022b]
	Cần chọn $3$ người đi công tác từ một tổ  có $30$ người, khi đó số cách chọn là
	 \choice
		{$\mathrm{A}_{30}^3$} 
		{$3^{30}$} 
		 {$10$}
		{\True $\mathrm{C}_{30}^3$} 
	 \loigiai{
		Số cách chọn $3$ người bất  kì trong $30$ người là $\mathrm{C}_{30}^3$.
	} 
\end{ex}

% nguồn: MaHoaNH_TracNghiem_Lop10
\begin{ex}%[0D8N2-2]%[ma:d9502c66]
	Có bao nhiêu khả năng có thể xảy ra đối với thứ tự giữa các đội trong một giải bóng có $4$ đội bóng tham gia (giả sử rằng không có hai  đội nào có điểm trùng nhau)?
	 \choice
		 {$4$}
		 {$10$}
		 {$16$}
		{\True $24$} 
	\loigiai{ 
		Mỗi kết quả thứ tự của $4$ đội bóng là một hoán vị của $4$ phần  tử.\\
	 Vậy có tất cả $4!=24$ khả năng xảy ra.}
\end{ex}

% nguồn: MaHoaNH_TracNghiem_Lop10
\begin{ex}%[0D8N2-2]%[ma:a90d89a7]
	Số các tổ hợp chập $k$ của $n$ phần  tử được tính bởi công thức
	\choice 
		{$\mathrm{C}_{n}^{k} =  \dfrac{n!}{k!}$}
		{$\mathrm{C}_{n}^{k} = \dfrac{n}{k(n-k)}$} 
		{\True $\mathrm{C}_{n}^{k}  = \dfrac{n!}{k!(n-k)!}$}
		{$\mathrm{C}_{n}^{k}  = \dfrac{n!}{(n-k)!}$}
	\loigiai{ 
		Số các tổ hợp chập  $k$ của $n$ phần tử được tính bởi công thức $\mathrm{C}_{n}^{k} = \dfrac{n!}{k!(n-k)!}$.
	 }
\end{ex}

% nguồn: MaHoaNH_TracNghiem_Lop10
\begin{ex}%[0D8N2-2]%[ma:46abdea1]
	Số các chỉnh hợp chập $k$ của $n$ phần tử được kí hiệu là 
	 \choice
		{$P_{n}$} 
		{$k!$} 
		 {\True $\mathrm{A}_{n}^{k}$}
		{$\mathrm{C}_{n}^{k}$} 
	 \loigiai{
		Số các chỉnh hợp chập $k$ của $n$ phần  tử được kí hiệu là $\mathrm{A}_{n}^{k}$.
	} 
\end{ex}

% nguồn: MaHoaNH_TracNghiem_Lop10
\begin{ex}%[0D8N2-2]%[ma:5865ac2f]
	Cho tập hợp $X$ gồm $10$ phần tử. Số các hoán vị  của $10$ phần tử của tập hợp $X$ là
	\choice 
		 {$2^{10}$}
		 {$10^{10}$}
		{\True$10!$} 
		 {$10^2$}
	 \loigiai{
		Số các hoán vị  của $10$ phần tử của tập hợp $X$ là $10!$.
	 }
\end{ex}

% nguồn: MaHoaNH_TracNghiem_Lop10
\begin{ex}%[0D8N2-2]%[ma:baca6487]
	Số cách xếp $6$ học sinh vào $6$ trong $10$ ghế trên một hàng  ngang sao cho mỗi học sinh ngồi một ghế là
	\choice 
		 {\True $\mathrm{A}_{10}^6$}
		{$\mathrm{C}_{10}^4$} 
		 {$6!$}
		{$6^{10}$} 
	\loigiai{ 
		Mỗi cách sắp xếp $6$ học sinh vào $6$ trong $10$ ghế trên một hàng ngang sao cho mỗi học sinh ngồi một  ghế là một chỉnh hợp chập $6$ của $10$ nên số cách xếp là $\mathrm{A}_{10}^6$.
	 }
\end{ex}

% nguồn: MaHoaNH_TracNghiem_Lop10
\begin{ex}%[0D8N2-2]%[ma:26c24b39]
	Số cách chọn $3$ học sinh từ $5$ học  sinh là
	 \choice
		 {$3!$}
		{\True  $\mathrm{C}_5^3$}
		 {$15$}
		 {$\mathrm{A}_5^3$}
	\loigiai{ 
		Số cách chọn $3$ học sinh từ $5$  học sinh là $\mathrm{C}_5^3$.
	 }
\end{ex}

% nguồn: MaHoaNH_TracNghiem_Lop10
\begin{ex}%[0D8N2-2]%[ma:9ae9a3bb]
	Có bao nhiêu cách xếp $5$ học  sinh thành một hàng dọc?
	 \choice
		{$5$} 
		{\True $5!$} 
		 {$4!$}
		 {$5^5$}
	\loigiai{ 
		Số cách sắp xếp $5$  học sinh vào $5$ vị trí là $5!$ cách xếp.
	 }
\end{ex}

% nguồn: MaHoaNH_TracNghiem_Lop10
\begin{ex}%[0D8H2-2]%[ma:b8dcb7b4]
	Nếu $\mathrm{C}_n^3=10$ thì $n$ có giá  trị là
	 \choice
		{$6$} 
		{\True  $5$}
		 {$3$}
		 {$4$}
	\loigiai{Điều  kiện $n\in\mathbb{N}$, $n\geq 3$.\\
	Ta  có $\mathrm{C}_n^3=10\Leftrightarrow\dfrac{n(n-1)(n-2)}{6}=10\Leftrightarrow n^3-3n^2+2n-60=0\Leftrightarrow n=5$.}
\end{ex}

% nguồn: MaHoaNH_TracNghiem_Lop10
\begin{ex}%[0D8H2-2]%[ma:307f0c8a]
	Trong một tuần bạn $A$ dự định mỗi ngày đi thăm một người bạn trong $12$ người  bạn của mình. Hỏi bạn $A$ có thể lập được bao nhiêu kế hoạch đi thăm bạn của mình (thăm một bạn không quá một lần).
	\choice 
		 {\True $3\,991\,680$}
		{$12!$} 
		{$35\,831\,808$} 
		 {$7!$}
	 \loigiai{
		Vì $1$ tuần có $7$ ngày nên  có $\mathrm{A}_{12}^7=3\,991\,680$ (kế hoạch).
	 }
\end{ex}

% nguồn: MaHoaNH_TracNghiem_Lop10
\begin{ex}%[0D8H2-2]%[ma:32fa7e7e]
	Tổ $1$ của lớp $10 \mathrm{A}_1$ có $10$ học sinh. Giáo viên cần chọn ra $3$ học sinh của tổ $1$ để tham dự trại hè. Hỏi giáo viên có bao nhiêu cách chọn? 
	 \choice
		 {\True $120$}
		 {$720$}
		{$27$} 
		 {$240$}
	 \loigiai{
		Số cách chọn $3$ học sinh từ $10$ học sinh của tổ $1$ là số tổ  hợp chập $3$ của $10$ phần tử.\\
	Vậy có  $\mathrm{C}_{10}^3=120$ cách.}
\end{ex}

% nguồn: MaHoaNH_TracNghiem_Lop10
\begin{ex}%[0D8H2-2]%[ma:cd85f257]
	Có bao nhiêu cách xếp $5$ quyển sách Văn khác nhau và $7$  quyển sách Toán khác nhau trên một kệ sách dài nếu các quyển sách Văn phải xếp kề nhau?
	 \choice
		 {$12!$}
		{$2\cdot  5!\cdot 7!$}
		 {\True $8!\cdot 5!$}
		{$5!.7!$} 
	 \loigiai{
		Ta coi $5$ quyển sách Văn là một Quyển và xếp Quyển này với $7$ quyển sách Toán khác nhau ta có $8!$ cách xếp. Mỗi  cách đổi vị trí các quyển sách văn cho nhau thì tương ứng sinh ra một cách xếp mới, mà có $5!$ cách đổi vị trí các quyển sách Văn.\\
		Vậy số cách xếp là $8!\cdot 5!$. 
	 }
\end{ex}

% nguồn: MaHoaNH_TracNghiem_Lop10
\begin{ex}%[0D8H2-2]%[ma:a4e9782b]
	Cho $n \in \mathbb{N}$,  $n \geq 4$ và $2\mathrm{C}_n^4=140$. Giá trị của $n$ bằng
	\choice 
		{$6$} 
		 {$9$}
		 {\True $8$}
		 {$7$}
	\loigiai{ 
		Ta có $2\mathrm{C}_n^4=140  \Rightarrow \mathrm{C}_n^4=70 \Rightarrow \dfrac{n!}{4!(n-4)!}=70$.\\
		Suy ra $n(n-1)(n-2)(n-3)=1680=8  \cdot 21 \cdot 10=8 \cdot 7 \cdot 3 \cdot 5 \cdot 2=8 \cdot 7 \cdot 6 \cdot 5 \Rightarrow n=8$.
	} 
\end{ex}

% nguồn: MaHoaNH_TracNghiem_Lop10
\begin{ex}%[0D8H2-2]%[ma:29a67896]
	Nghiệm của phương trình $\mathrm{A}_n^3=20n$  là
	 \choice
		 {$n=5$}
		 {\True $n=6$}
		{$n=8$} 
		{$n=9$} 
	 \loigiai{
		Điều kiện $\heva{&n\in \mathbb{N}  \\&n\ge 3.}$\\
		Ta có $\mathrm{A}_n^3=20n\Leftrightarrow n\left(n-1\right)\left(n-2\right)=20n\Leftrightarrow n^2-3n-18=0\Leftrightarrow \hoac{&n=6&&(\text{nhận})  \\&n=-3&&(\text{loại}).}$\\
		Vậy nghiệm của phương trình đã cho là  $n=6$.
	} 
\end{ex}

% nguồn: MaHoaNH_TracNghiem_Lop10
\begin{ex}%[0D8H2-2]%[ma:2bd4c950]
	Gọi $n_1$, $n_2$ là hai nghiệm của phương trình $\dfrac{1}{\mathrm{C}_n^1}-\dfrac{1}{\mathrm{C}_{n+1}^2}=\dfrac{7}{6\mathrm{C}_{n+4}^1}$.  Khi đó $n_1^2+n_2^2$ bằng
	 \choice
		{$85$} 
		{$64$} 
		 {$80$}
		{\True  $73$}
	\loigiai{ 
		Điều kiện $\heva{&n\geq 1\\ &n+1\geq 2\\ &n+4\geq 1 \\&n\in \mathbb{N}}\Leftrightarrow\heva{& n\geq 1  \\ & n\in\mathbb{N}.}$\\
		Khi đó \begin{eqnarray*}
		\dfrac{1}{\mathrm{C}_n^1}-\dfrac{1}{\mathrm{C}_{n+1}^2}=\dfrac{7}{6\mathrm{C}_{n+4}^1}  &
		\Leftrightarrow  & \dfrac{1}{n}-\dfrac{1}{\dfrac{(n+1)n}{2}}=\dfrac{7}{6(n+4)}\\
		&\Leftrightarrow & \dfrac{1}{n}-\dfrac{2}{n(n+1)}=\dfrac{7}{6(n+4)}\\ 
		&\Leftrightarrow  & 6(n+1)(n+4)-12(n+4)=7n(n+1)\\
		 &\Leftrightarrow & n^2-11n+24=0\\
		&\Leftrightarrow& \hoac{&n=8 &&(\text{nhận})  \\&n=3&&(\text{nhận}).}
	\end{eqnarray*}
Vậy $n_1^2+n_2^2=8^2+3^2=73$.} 
\end{ex}


%%%%%%%%%%%%%%%% ĐÚNG/SAI %%%%%%%%%%%%%%%%

% nguồn: MaHoaNH_DungSai_Lop10
\begin{ex}%[0D8H2-2]%[ma:0fe0d2de]
	Một lớp có $15$ nam  và $20$ nữ. Giáo viên chủ nhiệm cần chọn $4$ học sinh đi trực nhật. Trong các khẳng định sau, khẳng định nào đúng, khẳng định nào sai?
	 \choiceTF
		{\True Chọn $1$ nam và  $3$ nữ có $\mathrm{C}_{15}^1\cdot \mathrm{C}_{20}^3$ cách}
		{\True  Chọn $2$ nam, $2$ nữ có $\mathrm{C}_{15}^2\cdot \mathrm{C}_{20}^2$ cách}
		{Chọn $3$ nam,  $1$ nữ có $\mathrm{C}_{15}^3\cdot \mathrm{C}_{20}^1$ cách}
		{\True Số cách  chọn $4$ học sinh được chọn có cả nam và nữ là $\mathrm{C}_{35}^4$ cách}
	\loigiai{ 
		\begin{itemchoice}
			\itemch Đây là khẳng định đúng  vì \\
			Chọn $1$ nam, $3$ nữ có $\mathrm{C}_{15}^1\cdot  \mathrm{C}_{20}^3$ cách.
			\itemch Đây là  khẳng định đúng vì \\
			Chọn $2$ nam, $2$ nữ có $\mathrm{C}_{15}^2\cdot \mathrm{C}_{20}^2$ cách. 
			\itemch Đây là khẳng định đúng vì \\ 
			Chọn  $3$ nam, $1$ nữ có $\mathrm{C}_{15}^3\cdot \mathrm{C}_{20}^1$ cách.
			\itemch  Đây là khẳng định sai vì \\
			Mỗi cách chọn ra $4$ học sinh sao cho trong $4$ học sinh được chọn có cả nam và nữ. 
			\begin{itemize}
				\item \textbf{Trường  hợp 1}: Chọn $1$ nam, $3$ nữ có $\mathrm{C}_{15}^1\cdot \mathrm{C}_{20}^3$ cách.
				\item \textbf{Trường hợp 2}: Chọn $2$ nam, $2$ nữ có  $\mathrm{C}_{15}^2\cdot \mathrm{C}_{20}^2$ cách.
				\item \textbf{Trường hợp 3}: Chọn $3$ nam, $1$ nữ có $\mathrm{C}_{15}^3\cdot  \mathrm{C}_{20}^1$ cách.
			\end{itemize}
			Theo quy tắc cộng, số cách cách chọn $4$  trong đó có cả nam và nữ là $$\mathrm{C}_{15}^1\cdot \mathrm{C}_{20}^3+\mathrm{C}_{15}^2\cdot \mathrm{C}_{20}^2+\mathrm{C}_{15}^3\cdot \mathrm{C}_{20}^1=46150.$$
		\end{itemchoice}
	 }
\end{ex}

% nguồn: MaHoaNH_DungSai_Lop10
\begin{ex}%[0D8V2-2]%[ma:fec573e0]
	Từ $5$ bông hồng vàng, $3$ bông hồng trắng, $4$ bông hồng đỏ người ta chọn ra một  số bông hoa.
	 \choiceTF
		{\True Có $12$  cách chọn một bông hoa bất kỳ}
		{Có $20$ cách chọn một bông hồng  vàng và bông hồng trắng}
		{Có $60$ cách chọn bó hoa gồm  2 bông hồng đỏ và 1 bông hồng vàng}
		{ Có $130$ cách chọn sao cho có  ít nhất $3$ bông hồng vàng và $3$ bông hồng đỏ}
	\loigiai{ 
	\begin{itemchoice}
		\itemch Theo quy tắc cộng, có $5+3+4=12$ cách chọn một bông  hoa bất kỳ.
		\itemch Theo quy tắc  nhân, có $5\cdot 3=15$ cách chọn một bông hồng vàng và bông hồng trắng.
		\itemch Có $\mathrm{C}_4^2=6$ cách  chọn 2 bông hồng đỏ.\\
		Có $5$ cách chọn 2 bông hồng vàng.\\ 
		Theo quy tắc nhân, có $6\cdot 5=30$ cách chọn bó hoa gồm 2 bông  hồng đỏ và 1 bông hồng vàng.
		\itemch \textbf{Trường hợp 1:} chọn  $3$ bông hồng vàng, $3$ bông hồng đỏ và $1$ bông hồng trắng có $\mathrm{C}_5^3\cdot \mathrm{C}_4^3\cdot \mathrm{C}_3^1$ cách.\\
		\textbf{Trường hợp 2:} chọn $3$ bông hồng vàng, $4$ bông  hồng đỏ có $\mathrm{C}_5^3\cdot \mathrm{C}_4^4$ cách.\\
		\textbf{Trường hợp 3:} chọn $4$  bông hồng vàng, $3$ bông hồng đỏ có $\mathrm{C}_5^4\cdot \mathrm{C}_4^3$ cách.\\
		Vậy có $\mathrm{C}_5^3\cdot \mathrm{C}_4^3\cdot \mathrm{C}_3^1 + \mathrm{C}_5^3\cdot \mathrm{C}_4^4+\mathrm{C}_5^4\cdot \mathrm{C}_4^3=150$ cách chọn $7$ bông  hoa thỏa mãn yêu cầu.
	\end{itemchoice}
	} 
\end{ex}

% nguồn: MaHoaNH_DungSai_Lop10
\begin{ex}%[0D8V2-2]%[ma:609cd64f]
	Cho phương trình $\mathrm{A}_x^3+\mathrm{C}_x^{x-3}=14 x$. Các mệnh đề sau đúng hay  sai?
	 \choiceTF
		{\True Điều kiện: $x \in \mathbb{N}$ và $x \geq 3$} 
		{Phương trình có  chung tập nghiệm với phương trình $x^2-3x-10=0$}
		{Phương  trình có $2$ nghiệm phân biệt}
		{\True Nghiệm của phương trình là số nguyên  tố}
	\loigiai{ 
		Điều kiện: $x \in  \mathbb{N}$ và $x \geq 3$.\\
		Khi đó \begin{eqnarray*}
		\mathrm{A}_x^3+\mathrm{C}_x^{x-3}=14x& \Leftrightarrow &  x(x-1)(x-2)+\dfrac{x(x-1)(x-2)}{6}=14x\\
		 &\Leftrightarrow & 7(x-1)(x-2)=84\quad
		 (\text{vì } x>0)\\
		&\Leftrightarrow & x^2-3x-10=0\\ 
		 &\Leftrightarrow & \hoac{&x=5&& (\text{nhận})\\&x=-2 && (\text{loại}).}
	\end{eqnarray*}
	Vậy phương trình có  nghiệm là $x=5$.
	\begin{itemchoice}
		\itemch Đúng. 
		\itemch Sai. Phương trình $x^2-3x-10=0$ có tập  nghiệm là $S=\{-2;5\}$.
		\itemch Sai.  Phương trình đã cho có đúng một nghiệm.
		\itemch Đúng. Nghiệm  $x=5$ là số nguyên tố.
	\end{itemchoice}
} 
\end{ex}

% nguồn: MaHoaNH_DungSai_Lop10
\begin{ex}%[0D8T2-2]%[ma:839d17ba]
	Cho $n\in \mathbb{N^*}, n\ge 5$ thỏa mãn $\mathrm{C}_n^5=2002$.  Chọn khẳng định đúng
	 \choiceTF
		{$\mathrm{A}_n^5=\mathrm{C}_n^5\cdot  5$}
		{$\mathrm{A}_n^5=\mathrm{C}_n^5$} 
		 {\True $\dfrac{n!}{(n-5)!5!}=2002$}
		 {\True $\mathrm{A}_n^5=240240$}
	\loigiai{ 
		Ta  có
		\begin{itemize}
			\item $\mathrm{A}_n^5=\mathrm{C}_n^5\cdot 5!=240240$. 
			 \item $\mathrm{C}_n^5=2002\Leftrightarrow \dfrac{n!}{(n-5)!5!}=2002.$
		\end{itemize}
	 }
\end{ex}

% nguồn: MaHoaNH_DungSai_Lop10
\begin{ex}%[0D8T2-2]%[ma:0513743e]
	Cho số  nguyên dương $n$ thỏa mãn đẳng thức sau $\mathrm{C}_n^3+\mathrm{A}_n^2=376-2n$. Khẳng định nào sau đây đúng?
	\choiceTF 
		{$n<11$} 
		 {$1<n<12$}
		 {\True $1<n<14$}
		{\True $n>10$} 
	\loigiai{ 
		Điều kiện  $n\in \mathbb{N^*}, n\ge 3$. Ta có
		\begin{eqnarray*}
			& &  \mathrm{C}_n^3+\mathrm{A}_n^2=376-2n\\
			&\Leftrightarrow &\dfrac{n!}{3!(n-3)!}+\dfrac{n!}{(n-2)!}=376-2n \\ 
			&\Leftrightarrow &  n(n-1)(n-2)+6n(n-1)=2256-12\\
			&\Leftrightarrow  & n^3+3n^2+8n-2256=0\\
			&\Leftrightarrow & n=12. 
		\end{eqnarray*}
	} 
\end{ex}


%%%%%%%%%%%%%%%% TRẢ LỜI NGẮN %%%%%%%%%%%%%%%%

% nguồn: MaHoaNH_TraLoiNgan_Lop10
\begin{bt}%[0D8H2-2]%[ma:a30ce7fb]
	Cho  $\mathrm{C}_n^{n-3}=1140$. Tính $A=\dfrac{\mathrm{A}_n^6+\mathrm{A}_n^5}{\mathrm{A}_n^4}$.
	\shortans{256} 
	 \loigiai{
		 Điều kiện$:\heva{&n\in \mathbb{N}\\&n\ge 6.}$\\
		 Ta có $\mathrm{C}_n^{n-3}=1140 \Leftrightarrow \dfrac{n(n-1)(n-2)}{6}=1140 \Leftrightarrow n^3-3n^2+2n-6840=0\Leftrightarrow n=20$.\\
	Vậy  $A=\dfrac{\mathrm{A}_{20}^6+\mathrm{A}_{20}^5}{\mathrm{A}_{20}^4}=256$.}
\end{bt}

% nguồn: MaHoaNH_TraLoiNgan_Lop10
\begin{bt}%[0D8H2-2]%[ma:759aa51f]
	Một biển số xe có mã $71-$C$4$  ở phần trên, phần dưới là một dãy gồm $5$ chữ số chọn từ các số $0$ đến $9$. Hỏi có thể tạo ra được bao nhiêu biển số xe thỏa điều kiện: Trong dãy $5$ chữ số ở phần dưới, hai chữ số cuối phải là $99$ và ba chữ số đầu phải khác nhau?
	 \shortans{720}
	\loigiai{ 
		Gọi dãy $5$ số ở phần  dưới dạng $\overline{a_1a_2a_399}$. Điều kiện: $a_i\neq a_j$ với $i\leq j$ và $i$, $j \in \{1;2;3\}$.\\
		Mỗi cách chọn $a_1$, $a_2$, $a_3$ là một chỉnh hợp $3$ phần tử của $10$. Do  đó, số cách chọn là $\mathrm{A}_{10}^3$.\\
		 Vậy số biển số xe thỏa điều kiện đã cho là $\mathrm{A}_{10}^3 = 720$.
	} 
\end{bt}

% nguồn: MaHoaNH_TraLoiNgan_Lop10
\begin{bt}%[0D8V2-2]%[ma:a0df6cb7]
	Từ các chữ số $0, 1, 2, 3, 4, 5,  6$ có thể lập được bao nhiêu số tự nhiên nhỏ hơn $1000$?
	 \shortans{343}
	\loigiai{ 
		Số có một chữ số  được lập thành từ các số $0, 1, 2, 3, 4, 5, 6$ là $7$ số.\\
		Gọi số có hai chữ số  cần lập có dạng $\overline{ab}$ khi đó $a\ne 0\Rightarrow a$ có $6$ cách chọn, $b$ có $7$ cách chọn. Do đó có $6\cdot 7=42$ số có hai chữ số.\\
		 Gọi số có ba chữ số cần lập có dạng $\overline{abc}$ khi đó $a\ne 0\Rightarrow a$ có $6$ cách chọn, $b$ có $7$ cách chọn và $c$ có $7$ cách chọn. Do đó có $6\cdot 7 \cdot 7 = 294$ số có ba chữ số.\\
		Vậy có $7+42+294=343$ số tự  nhiên nhỏ hơn $1000$.
	} 
\end{bt}

% nguồn: MaHoaNH_TraLoiNgan_Lop10
\begin{bt}%[0D8V2-2]%[ma:6af2c63d]
	Từ tập $X=\{1,  2, 3, 4, 5, 6, 7, 8, 9\}$ lập được bao nhiêu số tự nhiên có $4$ chữ số đôi một khác nhau sao cho không có hai chữ số liên tiếp nào cùng chẵn?
	 \shortans{1800}
	 \loigiai{
		Có các trường hợp sau thỏa mãn số lập được không có hai chữ  số liên tiếp nào cùng chẵn:
		\begin{itemize}
			\item Tất cả $4$ chữ số đều là số lẻ. Trường hợp này ta chọn $4$ trong $5$ chữ số lẻ của tập $X$ và sắp thứ tự nên  có $A^4_5$ số thỏa mãn.
			\item Có $3$ chữ số lẻ và đúng $1$ chữ số chẵn. Trường hợp  này ta chọn $3$ trong $5$ chữ số lẻ của tập $X$, thêm vào $1$ chữ số chẵn trong $4$ chữ số chẵn của tập $X$, sau đó hoán vị $4$ chữ số này. Như vậy có $C^3_5\cdot 4\cdot 4!$ số thỏa mãn.
			\item Có đúng $2$ chữ số lẻ và $2$ chữ số chẵn. Trường hợp này sắp xếp theo thứ tự L C L C hoặc C L C L hoặc C L L  C (với C là kí hiệu cho chữ số chẵn và L là kí hiệu cho chữ số lẻ). Như vậy có $5\cdot 4\cdot 4\cdot 3\cdot 3$ số thỏa mãn.
		\end{itemize}
		Vậy có tất cả $A^4_5+C^3_5\cdot 4\cdot 4!+5\cdot 4\cdot 4\cdot 3\cdot 3=1800$ số thỏa  mãn bài ra.
	 }
\end{bt}

% nguồn: MaHoaNH_TraLoiNgan_Lop10
\begin{bt}%[0D8V2-2]%[ma:27f55cb0]
	Từ $2$ chữ số $1$ và $8$, có thể lập được bao nhiêu số tự nhiên có $8$  chữ số sao cho không có hai chữ số $1$ đứng cạnh nhau?
	\shortans{55} 
	\loigiai{ 
	Gọi số tự nhiên $x$ cần tìm có dạng $x=\overline{abcdefgh}$; với $a$, $b$, $c$, $d$,  $e$, $f$, $g$, $h\in\left\{1;8\right\}$.\\
	Vì $x$ không có hai chữ số $1$ đứng cạnh nhau nên  ta có có trường hợp sau:
	\begin{itemize}
		\item $x$ chỉ gồm các chữ số $8$: Có  $1$ số là $88.888.888$.
		\item $x$ chỉ gồm $1$ chữ số $1$ và $7$ chữ số $8$: Có $8$ vị trí cho chữ số $1$,  tức là có $8$ số thỏa mãn.
		\item $x$ gồm $2$ chữ số số $1$ và $6$ chữ số $8$: Giữa $6$ chữ số $8$ ta có $5$ khoảng trống, cộng thêm với $2$ khoảng trống ở ngoài cùng bên trái và ngoài cùng bên phải thì tổng cộng ta có $7$ khoảng trống có thể đặt chữ số  $1$ vào. Chọn ra $2$ vị trí trong $7$ khoảng trống này, ta có $\mathrm{C}_7^2 = 21$ số.
		\item $x$ gồm $3$ chữ số số $1$ và $5$ chữ số $8$: Giữa $5$ chữ số $8$ ta có $4$ khoảng trống, cộng thêm với $2$ khoảng trống ở ngoài cùng bên trái và ngoài cùng bên phải thì tổng cộng ta có $6$ khoảng trống có thể đặt chữ số $1$ vào. Chọn ra $3$ vị trí trong $6$ khoảng trống này, ta có $\mathrm{C}_6^3 =  20$ số.
		\item $x$ gồm $4$ chữ số số $1$ và $4$ chữ số $8$: Giữa $4$ chữ số $8$ ta có $3$ khoảng trống, cộng thêm với $2$ khoảng trống ở ngoài cùng  bên trái và ngoài cùng bên phải thì tổng cộng ta có $5$ khoảng trống có thể đặt chữ số $1$ vào. Chọn ra $4$ vị trí trong $5$ khoảng trống này, ta có $\mathrm{C}_5^4 = 5$ số.
	\end{itemize}
	Vậy tổng  cộng ta có $1+8+21+20+5 = 55$ số.
	} 
\end{bt}
